\documentclass[../../main.tex]{subfiles}
\begin{document}
\section{奇异同调}\label{section:ch18-singular-homology}

为了叙述 de Rham 定理，我们需要一点代数拓扑。本节只引入所需的最少概念。

设 $v_0,\dots,v_p$ 是 $\mathbb{R}^n$ 中一组仿射线性无关的点。它们的\emph{几何 $p$-单形}定义为所有凸组合
\begin{align*}
t_0v_0+\cdots+t_pv_p,\qquad t_i\geqslant0,\quad t_0+\cdots+t_p=1
\end{align*}
构成的集合，记为 $[v_0,\dots,v_p]$。每个 $[v_{i_0},\dots,v_{i_k}]$ 称为该单形的一个面。$(p-1)$ 维面称为它的\emph{边界面}。恰有 $p+1$ 个边界面，依次省略每个顶点得到；第 $i$ 个边界面 $[v_0,\dots,\widehat{v_i},\dots,v_p]$（省略 $v_i$）记为 $\partial_i[v_0,\dots,v_p]$，称为与 $v_i$ 相对的面。

\begin{exercise}\label{exercise:ch18-01}
证明一个几何 $p$-单形是一个 $p$ 维光滑带角流形，并光滑嵌入在 $\mathbb{R}^n$ 中。
\end{exercise}

\emph{标准 $p$-单形}是单形 $\Delta^p=[e_0,e_1,\dots,e_p]\subseteq\mathbb{R}^p$，其中 $e_0=0$，$e_i$ 是第 $i$ 个标准基向量。例如，$\Delta^0=\{0\}$，$\Delta^1=[0,1]$，$\Delta^2$ 是以 $(0,0),(1,0),(0,1)$ 为顶点的三角形连同其内部，$\Delta^3$ 是一个实心四面体（\cref{figure:ch18-standard-simplex}）。

\begin{figure}[ht]
\centering
\includegraphics[width=0.5\textwidth]{fig-18-1.png}
\caption{标准单形}\label{figure:ch18-standard-simplex}
\end{figure}

设 $M$ 是一个拓扑空间。一个连续映射 $\sigma:\Delta^p\to M$ 称为 $M$ 中的一个\emph{奇异 $p$-单形}。$M$ 的 $p$ 次\emph{奇异链群}，记为 $C_p(M)$，是由 $M$ 中所有奇异 $p$-单形生成的自由阿贝尔群。这个群的一个元素称为一个\emph{奇异 $p$-链}，是奇异 $p$-单形以整数为系数的有限形式线性组合。

一个常见的特殊情形是 $M$ 为某个欧氏空间 $\mathbb{R}^m$ 的凸子集。此时对任意有序 $(p+1)$ 元组 $(w_0,\dots,w_p)\in M$（不必仿射线性无关），存在唯一的仿射映射 $\mathbb{R}^p\to\mathbb{R}^m$ 把 $e_i$ 映到 $w_i$。该仿射映射在 $\Delta^p$ 上的限制记为 $A(w_0,\dots,w_p)$，称为 $M$ 中的一个\emph{仿射奇异单形}。

\begin{figure}[ht]
\centering
\includegraphics[width=0.55\textwidth]{fig-18-2.png}
\caption{奇异边界算子}\label{figure:ch18-singular-boundary}
\end{figure}

对每个 $i=0,\dots,p$，定义 $\Delta^p$ 中的第 $i$ 个\emph{面映射} $F_{i,p}:\Delta^{p-1}\to\Delta^p$ 为仿射奇异 $(p-1)$-单形
\begin{align*}
F_{i,p}=A(e_0,\dots,\widehat{e_i},\dots,e_p).
\end{align*}
它把 $\Delta^{p-1}$ 同胚地映到边界面 $\partial_i\Delta^p$ 上。

一个奇异 $p$-单形 $\sigma:\Delta^p\to M$ 的\emph{边界}是奇异 $(p-1)$-链
\begin{align}
\partial\sigma=\sum_{i=0}^p(-1)^i\sigma\circ F_{i,p}.
\end{align}
例如，若 $\sigma$ 是一个奇异 $2$-单形，其边界是三个奇异 $1$-单形以 $\pm1$ 为系数的形式和（\cref{figure:ch18-singular-boundary}）。这唯一地延拓为群同态 $\partial:C_p(M)\to C_{p-1}(M)$，称为\emph{奇异边界算子}。

\begin{lemma}\label{lemma:ch18-singular-dd-zero}
若 $c$ 是任意奇异链，则 $\partial(\partial c)=0$。
\end{lemma}

\begin{proof}[证明梗概]
出发点是事实：当 $i>j$ 时
\begin{align*}
F_{i,p}\circ F_{j,p-1}=F_{j,p}\circ F_{i-1,p-1},
\end{align*}
这可以通过考察两个复合对 $\Delta^{p-2}$ 每个顶点的作用来验证。利用这一点，引理的证明就是一个直接的计算。
\end{proof}

一个奇异 $p$-链 $c$ 若满足 $\partial c=0$ 则称为一个\emph{闭链}；若 $c=\partial b$ 对某个奇异 $(p+1)$-链 $b$ 成立，则称为一个\emph{边缘}。记 $Z_p(M)$ 为 $M$ 中奇异 $p$-闭链的集合，$B_p(M)$ 为奇异 $p$-边缘的集合。由于 $\partial$ 是同态，$Z_p(M)$ 与 $B_p(M)$ 都是 $C_p(M)$ 的子群；又由于 $\partial\circ\partial=0$，它们满足 $B_p(M)\subseteq Z_p(M)$。$M$ 的第 $p$ 个\emph{奇异同调群}定义为商群
\begin{align*}
H_p(M)=Z_p(M)/B_p(M).
\end{align*}
换句话说，阿贝尔群与同态的序列
\begin{align*}
\cdots\to C_{p+1}(M)\xrightarrow{\partial}C_p(M)\xrightarrow{\partial}C_{p-1}(M)\to\cdots
\end{align*}
是一个复形，称为\emph{奇异链复形}，而 $H_p(M)$ 是该复形的第 $p$ 个同调群。一个奇异 $p$-闭链 $c$ 在 $H_p(M)$ 中的等价类称为它的\emph{同调类}，记为 $[c]$。若两个 $p$-闭链相差一个边缘，则称它们\emph{同调}。

\begin{figure}[ht]
\centering
\includegraphics[width=0.6\textwidth]{fig-18-3.png}
\caption{连续映射诱导的同调同态}\label{figure:ch18-induced-homology}
\end{figure}

一个连续映射 $F:M\to N$ 在每个奇异链群上诱导一个同态 $F_\#:C_p(M)\to C_p(N)$，定义为 $F_\#(\sigma)=F\circ\sigma$，并线性延拓到链上。直接计算表明 $F_\#\circ\partial=\partial\circ F_\#$，故 $F_\#$ 是一个链映射，从而在奇异同调群上诱导一个同态，记为 $F_*:H_p(M)\to H_p(N)$。显然 $(G\circ F)_*=G_*\circ F_*$，$(\operatorname{Id}_M)_*=\operatorname{Id}_{H_p(M)}$。因此第 $p$ 个奇异同调定义了从拓扑空间与连续映射范畴到阿贝尔群与同态范畴的一个协变函子。特别地，同胚的空间有同构的奇异同调群。

\begin{proposition}[奇异同调群的性质]\label{proposition:ch18-singular-homology-properties}
\begin{enumerate}[(a)]
\item 对任意单点空间 $\{q\}$，$H_0(\{q\})$ 是由映 $\Delta^0$ 到 $q$ 的唯一奇异 $0$-单形的同调类生成的无穷循环群，且 $H_p(\{q\})=0$ 对所有 $p\neq0$。
\item 设 $\{M_j\}$ 是任意一族拓扑空间，$M=\coprod_j M_j$。包含映射 $\iota_j:M_j\hookrightarrow M$ 诱导同构 $\bigoplus_j H_p(M_j)\cong H_p(M)$。
\item 同伦等价的空间有同构的奇异同调群。
\end{enumerate}
\end{proposition}

\begin{theorem}[奇异同调的 Mayer--Vietoris]\label{theorem:ch18-mayer-vietoris-singular}
设 $M$ 是一个拓扑空间，$U,V\subseteq M$ 是其并为 $M$ 的开子集。对每个 $p$，存在一个连接同态 $\partial_*:H_p(M)\to H_{p-1}(U\cap V)$，使下面的序列正合：
\begin{align*}
\cdots\to H_p(U\cap V)\xrightarrow{\alpha}H_p(U)\oplus H_p(V)\xrightarrow{\beta}H_p(M)\xrightarrow{\partial_*}H_{p-1}(U\cap V)\to\cdots.
\end{align*}
\end{theorem}
\end{document}
